\documentclass[11pt]{article}
\usepackage[margin=1in]{geometry}
\usepackage{amsmath, amssymb}
\usepackage{booktabs}
\usepackage{graphicx}
\usepackage{parskip}
\usepackage{microtype}
\usepackage{hyperref}
\hypersetup{colorlinks=true, linkcolor=black, urlcolor=blue}

\title{Pricing Project\\[0.4ex]
\large Choice Modeling, Assortment Optimization and Pricing}
\author{Alexander Dyckerhoff, Parsa Bazargani, 
Sam DeLucia-Green }
\date{April 2026}

\begin{document}

\maketitle
\tableofcontents
\newpage

%%===================================================================
\section{Problem 1: MNL Model}
%%===================================================================

\begin{table}[h]
\centering
\begin{tabular}{lr}
\toprule
Feature & Coefficient (standardized) \\
\midrule
Intercept                   & $-1.747$ \\
Star rating                 & $+0.408$ \\
Review score                & $+0.109$ \\
Brand (bool)                & $+0.101$ \\
Location score              & $+0.022$ \\
Accessibility score         & $+0.043$ \\
Log historical price        & $-0.067$ \\
Price (USD)                 & $-1.331$ \\
Promotion flag              & $+0.160$ \\
\bottomrule
\end{tabular}
\caption{Estimated MNL coefficients. Features are standardized; each coefficient is the utility change per one-std-dev increase.}
\end{table}

Both standardized and raw (original-unit) coefficients are computed. Standardized coefficients allow direct comparison of feature importance by effect size; raw coefficients give interpretable magnitudes (e.g., per-dollar utility change).

\textbf{Price ($-1.33$)} is the dominant driver: a one-std-dev price increase drops utility by 1.33 units, enough to flip most ranking decisions. \textbf{Star rating ($+0.41$)} is the strongest positive quality signal. \textbf{Promotion flag ($+0.16$)}, \textbf{review score ($+0.11$)}, \textbf{brand ($+0.10$)}, \textbf{accessibility ($+0.04$)}, and \textbf{location ($+0.02$)} are all positive but smaller, contributing incrementally to utility. \textbf{Log historical price ($-0.07$)} is a small penalty: hotels that were historically more expensive are slightly less attractive holding current price fixed. The \textbf{intercept ($-1.75$)} is negative, indicating a high baseline threshold: customers require a meaningful combination of positive features before any hotel becomes their preferred choice.

%%===================================================================
\section{Problem 2: Assortment Optimization under MNL}
%%===================================================================

\textbf{Setup.} Under MNL, expected revenue of assortment $S$ is
\[
R(S) = \frac{\sum_{j \in S} p_j v_j}{1 + \sum_{j \in S} v_j}.
\]
We use the Talluri--van Ryzin revenue-ordered assortment result: an optimal assortment has the form $S(r) = \{j : p_j \ge r\}$ for some threshold $r$. Sort by price descending, form every prefix, and pick the one with highest $R(S)$. This is $O(n \log n)$ and exact (no MILP needed under a single MNL).

\begin{table}[h]
\centering
\begin{tabular}{lrrr}
\toprule
Dataset & Hotels & Offered & Expected revenue \\
\midrule
data1 & 27 & 18 & \$107.35 \\
data2 & 29 & 10 & \$131.11 \\
data3 & 26 & 18 & \$121.05 \\
data4 & 27 & 11 & \$97.41  \\
\bottomrule
\end{tabular}
\caption{Optimal assortments under MNL.}
\end{table}

\textbf{Why some hotels are dropped.} Adding a low-price hotel raises the denominator ($1 + \sum v_j$) through its preference weight $v_j$, but contributes little to the numerator ($p_j v_j$) because $p_j$ is small. Expected revenue falls whenever the hotel's price is below the current ratio. The offered prices in each dataset skew heavily to the high end, which is exactly what the revenue-ordered property predicts.

\textbf{Why $|S|$ varies across datasets.} data2 and data4 keep far fewer hotels than data1 and data3. This is driven by the price-to-preference ratios within each dataset. When the top-priced hotels have large $v_j$, adding lower-priced items hurts more, so the optimal cutoff is stricter. data2's offered prices start at \$133 and stretch to \$233; data4's cutoff excludes everything below \$100. data1 and data3 have tighter price distributions, so more hotels clear the threshold.

%%===================================================================
\section{Problem 3: Pricing under MNL}
%%===================================================================

\subsection*{Equal-markup derivation}

Differentiating $R(S) = \sum_{j \in S} p_j v_j / (1 + \sum_{j \in S} v_j)$ with respect to $p_j$ and using $\partial v_j / \partial p_j = \beta_p v_j$ (where $\beta_p = \beta_{\text{price,std}} / \sigma_{\text{price}}$ is the raw-scale price coefficient) gives:
\[
\frac{\partial R}{\partial p_j} = 0 \implies p_j^* = R^*(S) + \frac{1}{|\beta_p|}.
\]
Since the right-hand side is identical for every $j$, all hotels must have the same optimal price: $p_1^* = \cdots = p_n^* = p^*$. This is the classical MNL equal-markup result. Numerically, $\beta_p = -1.331 / \sigma_{\text{price}}$. For data1: $p^* \approx \$314.35$ and $R^*(S^*) \approx \$177.80$, so $1/|\beta_p| \approx 136.55$, implying $\sigma_{\text{price}} \approx 181.7$, consistent with the wide price range in the data.

\begin{table}[h]
\centering
\begin{tabular}{lrrrr}
\toprule
Dataset & Current revenue & Optimal revenue & Lift & Optimal price \\
\midrule
data1 & \$104.65 & \$177.80 & $1.70\times$ & $\approx\$314$ \\
data2 & \$113.36 & \$248.80 & $2.19\times$ & $\approx\$385$ \\
data3 & \$118.17 & \$176.47 & $1.49\times$ & $\approx\$313$ \\
data4 & \$78.69  & \$214.47 & $2.73\times$ & $\approx\$351$ \\
\bottomrule
\end{tabular}
\caption{Pricing results under MNL. All hotels displayed; prices optimized jointly.}
\end{table}

Optimal prices collapse to a near-constant value per dataset. Hotel-specific features appear only in the revenue weights $v_j$, not in the optimal price itself. Revenue lifts range from $1.5\times$ (data3) to $2.7\times$ (data4) over status quo prices. All four optima sit well inside the bounds (approximately \$313--\$385 vs.\ upper bound $\geq \$1{,}000$), so the bounds are not binding; the optimum is interior.

The single-price result is unrealistic in practice because it assumes homogeneous price sensitivity across all hotels; relaxing that (e.g., per-hotel price coefficients) would yield differentiated prices.

%%===================================================================
\section{Problem 4: Mixture of MNL}
%%===================================================================

We estimate separate MNL models for two customer types defined by booking window:
\begin{itemize}
  \item \textbf{Type 1 (early):} booking window $\geq 7$ days (4,537 queries, 54.3\%)
  \item \textbf{Type 2 (late):} booking window $< 7$ days (3,817 queries, 45.7\%)
\end{itemize}

\begin{table}[h]
\centering
\begin{tabular}{lrrr}
\toprule
Feature & Type 1 (early) & Type 2 (late) & Diff (late $-$ early) \\
\midrule
Intercept            & $-1.918$ & $-1.540$ & $+0.378$ \\
Star rating          & $+0.379$ & $+0.463$ & $+0.084$ \\
Review score         & $+0.125$ & $+0.096$ & $-0.029$ \\
Brand                & $+0.092$ & $+0.113$ & $+0.021$ \\
Location score       & $-0.021$ & $+0.087$ & $+0.109$ \\
Accessibility        & $+0.058$ & $+0.026$ & $-0.032$ \\
Log hist.\ price     & $-0.095$ & $-0.032$ & $+0.063$ \\
Price (USD)          & $-1.072$ & $-1.697$ & $-0.625$ \\
Promotion flag       & $+0.135$ & $+0.194$ & $+0.059$ \\
\bottomrule
\end{tabular}
\caption{MMNL coefficients (standardized features) by customer type.}
\end{table}

\textbf{Price sensitivity.} Late bookers have a much more negative price coefficient ($-1.70$) than early bookers ($-1.07$): approximately 58\% more price-sensitive. A possible explanation is that late bookers are leisure travelers comparing last-minute deals, while early bookers have already committed to a trip and are less willing to walk away over price.

\textbf{Star rating.} Both types value star rating positively, but late bookers weight it more ($+0.46$ vs.\ $+0.38$). Combined with higher price sensitivity, late customers appear to want good quality at a good price.

\textbf{Promotion flag.} Both types respond positively to promotions, but the effect is stronger for late bookers ($+0.19$ vs.\ $+0.13$), consistent with late bookers being more deal-driven.

\textbf{Location score.} Early bookers are essentially indifferent to location ($-0.02$), while late bookers value it positively ($+0.09$). Customers booking for tonight care where the hotel is; customers booking months ahead have less concrete plans.

\textbf{Historical price.} Both coefficients are negative, but early bookers weight it more ($-0.09$ vs.\ $-0.03$): they penalize historically expensive hotels more strongly.

\textbf{Review score, brand, and accessibility.} Differences are below 0.04 in magnitude and do not suggest meaningful behavioral distinctions.

\textbf{Intercept.} The intercept is less negative for late bookers ($-1.54$ vs.\ $-1.92$), meaning late customers have a higher baseline probability of booking. This fits the intuition that customers searching close to their check-in date are more committed to booking.

\textbf{Summary.} Late bookers are more price-sensitive, more responsive to promotions and star rating, care more about location, and are more likely to book overall. Early bookers are less price-sensitive but more cautious about historical price, and have a lower baseline booking probability. Overall, late customers behave like deal-seeking, decision-ready buyers, while early customers behave like lower-intent browsers.

%%===================================================================
\section{Problem 5: Early vs.\ Late Reservations}
%%===================================================================

\subsection*{MILP formulation}

Under the Mixture of MNL, expected revenue of assortment $S$ is
\[
R(S) = \sum_k \theta_k \cdot \frac{\displaystyle\sum_j p_j v_{jk} x_j}{1 + \displaystyle\sum_j v_{jk} x_j},
\]
where $x_j \in \{0,1\}$ indicates whether hotel $j$ is in $S$. The ratio per type is non-linear in $x$, so we apply the Charnes--Cooper linearization.

For each type $k$, introduce $y_k = 1 / (1 + \sum_j v_{jk} x_j)$ and $w_{jk} = y_k x_j$. The MILP is:
\[
\max \sum_k \theta_k \sum_j p_j v_{jk} w_{jk}
\]
subject to:
\begin{itemize}
  \item \textbf{Normalization:} $y_k + \sum_j v_{jk} w_{jk} = 1$ for each type $k$
  \item \textbf{Linking (McCormick):} $w_{jk} \le x_j$,\; $w_{jk} \le y_k$,\; $w_{jk} \ge y_k - (1 - x_j)$,\; $w_{jk} \ge 0$
  \item $x_j \in \{0,1\}$,\; $y_k \in [0,1]$,\; $w_{jk} \in [0,1]$
\end{itemize}

When $x_j = 0$ the linking forces $w_{jk} = 0$. When $x_j = 1$ it forces $w_{jk} = y_k$. The normalization then implies $y_k = 1 / (1 + \sum_{j \in S} v_{jk})$, recovering the original ratio. The objective linearizes to $\sum_k \theta_k y_k (\sum_{j \in S} p_j v_{jk})$, which is exactly the mixture expected revenue.

For the type-conditional assortments $S_1$ and $S_2$ we use the revenue-ordered property (optimal under a single MNL), avoiding the MILP.

\subsection*{Results}

\begin{table}[h]
\centering
\resizebox{\textwidth}{!}{%
\begin{tabular}{lrrrrrrrrrr}
\toprule
Dataset & $|S|$ & $R_\text{mix}(S)$ & $|S_1|$ & $R_1(S_1)$ & $R_1(S)$ & gap$_1$ & $|S_2|$ & $R_2(S_2)$ & $R_2(S)$ & gap$_2$ \\
\midrule
data1 & 18 & 107.30 & 21 & 103.47 & 103.25 & 0.22 & 16 & 112.22 & 112.12 & 0.10 \\
data2 & 10 & 131.28 & 10 & 126.92 & 126.92 & 0.00 &  7 & 137.38 & 136.46 & 0.92 \\
data3 & 18 & 121.00 & 19 & 117.52 & 117.44 & 0.08 & 17 & 125.35 & 125.23 & 0.12 \\
data4 & 11 &  97.38 & 11 &  94.69 &  94.69 & 0.00 & 10 & 100.64 & 100.58 & 0.06 \\
\bottomrule
\end{tabular}%
}
\caption{Problem 5: compromise vs.\ type-specific assortments. $S$ is mixture-optimal; $S_1$, $S_2$ are optimal for early and late customers respectively. Revenues in \$.}
\end{table}

\textbf{The gaps are small.} The cost of not knowing the customer's type is at most \$0.92 (data2, type 2) and usually under \$0.25. The compromise assortment $S$ performs nearly as well as the type-specific assortments because the two types differ mainly in the magnitude of preferences, not the direction.

\textbf{$S_1$ is always larger than $S_2$.} In every dataset, the optimal assortment for late customers contains fewer hotels (21 vs.\ 16 in data1, 10 vs.\ 7 in data2, 19 vs.\ 17 in data3, 11 vs.\ 10 in data4). Late customers are more price-sensitive, so adding lower-priced hotels hurts revenue more for type 2.

\textbf{$S$ sits between $S_1$ and $S_2$ in size.} In data1 and data3, $|S|$ falls between $|S_1|$ and $|S_2|$, reflecting a weighted balance of both types' preferences.

\textbf{When types agree, the gap is zero.} In data2 and data4, $S = S_1$ exactly. Late customers absorb a small revenue loss (\$0.92 and \$0.06) under this shared assortment.

\textbf{Main takeaway.} Segmentation improves model fit, but on the assortment decision the value of knowing the customer type is small in this data. The value of personalization depends on how much segments disagree on the \emph{decision}, not just on the \emph{parameters}.

%%===================================================================
\section{Problem 6: Mixture of MNL with Other Customer Type Definitions}
%%===================================================================

\subsection*{Customer type definition}

We redefine customer types based on the number of adults in the search query:
\begin{itemize}
  \item \textbf{Type 1 (solo):} \texttt{srch\_adults\_count == 1} (2,130 queries, 25.5\%)
  \item \textbf{Type 2 (group):} \texttt{srch\_adults\_count > 1} (6,224 queries, 74.5\%)
\end{itemize}

Solo travelers are often business travelers on a single budget; group bookings (couples, families) split cost across adults and tend to be leisure-oriented. The Saturday/non-Saturday split suggested in the problem statement was considered but discarded: Saturday vs.\ non-Saturday produces near-identical MNL coefficients across all eight features (the full comparison is in Section~\ref{sec:sat}), meaning the split captures day-of-week noise rather than a genuine behavioral distinction. Solo/group was chosen because it encodes a structural difference in travel purpose that maps onto measurable preference heterogeneity -- particularly brand loyalty and location sensitivity, which are proxies for business vs.\ leisure intent.

\subsection*{Parameter analysis}

\begin{table}[h]
\centering
\begin{tabular}{lrrr}
\toprule
Feature & Solo & Group & Diff (group $-$ solo) \\
\midrule
Intercept            & $-1.365$ & $-1.870$ & $-0.505$ \\
Star rating          & $+0.435$ & $+0.392$ & $-0.043$ \\
Review score         & $+0.073$ & $+0.125$ & $+0.052$ \\
Brand                & $+0.158$ & $+0.081$ & $-0.077$ \\
Location score       & $+0.081$ & $-0.001$ & $-0.082$ \\
Accessibility        & $+0.024$ & $+0.050$ & $+0.027$ \\
Log hist.\ price     & $-0.060$ & $-0.068$ & $-0.008$ \\
Price (USD)          & $-1.440$ & $-1.283$ & $+0.157$ \\
Promotion flag       & $+0.141$ & $+0.169$ & $+0.028$ \\
\bottomrule
\end{tabular}
\caption{MMNL coefficients (standardized features) for solo vs.\ group split.}
\end{table}

Solo travelers are more price-sensitive ($-1.44$ vs.\ $-1.28$) despite paying alone; group bookings split cost across adults, lowering per-person sensitivity. Solo travelers weight brand and location roughly $2\times$ more heavily, the signature of business travel where loyalty programs and proximity are constraints. Group travelers lean more on review score and accessibility. Compared to Problem 4's early/late split (price gap: 0.63 units), this split produces a smaller price gap (0.16 units) but a larger intercept gap (0.50) and more differentiation on brand and location: early/late separates on readiness to buy, solo/group on travel purpose.

\subsection*{Assortment results}

\begin{table}[h]
\centering
\resizebox{\textwidth}{!}{%
\begin{tabular}{lrrrrrrrrrr}
\toprule
Dataset & $|S|$ & $R_\text{mix}(S)$ & $|S_s|$ & $R_s(S_s)$ & $R_s(S)$ & gap$_s$ & $|S_g|$ & $R_g(S_g)$ & $R_g(S)$ & gap$_g$ \\
\midrule
data1 & 18 & 107.23 & 16 & 114.94 & 114.55 & 0.39 & 21 & 104.84 & 104.72 & 0.12 \\
data2 & 10 & 130.39 &  6 & 144.40 & 141.03 & 3.36 & 10 & 126.74 & 126.74 & 0.00 \\
data3 & 18 & 120.78 & 15 & 132.40 & 131.64 & 0.76 & 19 & 117.17 & 117.07 & 0.10 \\
data4 & 11 &  97.26 & 10 & 102.29 & 102.06 & 0.24 & 11 &  95.62 &  95.62 & 0.00 \\
\bottomrule
\end{tabular}%
}
\caption{Problem 6: compromise vs.\ type-specific assortments. $S_s$ is optimal for solo, $S_g$ for group. Revenues in \$.}
\end{table}

\textbf{Gaps are larger than in Problem 5.} The cost of the compromise assortment for solo customers rises to \$3.36 on data2, averaging above \$0.75 in three of four datasets. Group-side gaps are very small (0.00--0.12). The mixture is pulled toward the majority (group) type, so the minority (solo) absorbs most of the lost revenue.

\textbf{$|S_s|$ is always smaller than $|S_g|$.} Solo customers are offered fewer hotels in every dataset (16 vs.\ 21 in data1, 6 vs.\ 10 in data2, 15 vs.\ 19 in data3, 10 vs.\ 11 in data4). A few hotels dominate solo preferences, so the optimal assortment concentrates on those top items.

\textbf{$S$ coincides with $S_g$ in data2 and data4.} Because group customers are 74.5\% of the population, the mixture assortment inherits their preferred assortment on these two datasets, and solo customers absorb the full adaptation cost.

\textbf{Contrast with Problems 5 and Appendix~\ref{sec:sat}.} Problem 5's early/late split produced gaps up to \$0.92. The Saturday/non-Saturday split (Appendix~\ref{sec:sat}) produces gaps of at most \$0.01. This solo/group split produces a gap of \$3.36 on data2, more than $3\times$ larger than early/late and $300\times$ larger than Saturday. The ordering -- solo/group $\gg$ early/late $\gg$ Saturday -- confirms that trip-purpose signals drive real assortment-level lift, while demographic-timing signals produce minimal value of personalization.

\textbf{Main takeaway.} Solo vs.\ group disagrees on the assortment itself: solo wants fewer, higher-priced, brand-heavy hotels; group wants a broader assortment. That disagreement at the decision level is why personalization pays off here.

%%===================================================================
\section{Problem 7: AI Agents as Customers}
%%===================================================================

\subsection*{Coefficient comparison: human vs.\ AI agent}

\begin{table}[h]
\centering
\begin{tabular}{lrr}
\toprule
Feature & Human MNL & Agent MNL \\
\midrule
Intercept            & $-1.747$ & $-5.778$ \\
Star rating          & $+0.408$ & $+0.315$ \\
Review score         & $+0.109$ & $+0.718$ \\
Brand                & $+0.101$ & $+0.157$ \\
Location score       & $+0.022$ & $+0.074$ \\
Accessibility        & $+0.043$ & $+0.103$ \\
Log hist.\ price     & $-0.067$ & $+0.001$ \\
Price (USD)          & $-1.331$ & $-0.981$ \\
Promotion flag       & $+0.160$ & $+0.530$ \\
\bottomrule
\end{tabular}
\caption{MNL re-estimated on human data (Problem 1) vs.\ AI-generated choices (Problem 7).}
\end{table}

Both models agree on sign for 7 of 8 features, confirming the agent captures the broad structure of human preferences. The agent's intercept is far more negative ($-5.78$ vs.\ $-1.75$), meaning it chose the no-purchase option much more often than real customers. The agent overweights review score ($\approx 6.6\times$ larger) and promotion flag ($\approx 3.3\times$ larger) relative to humans, suggesting it reasons more ``rationally'' about quality signals. Price sensitivity is lower for the agent ($-0.98$ vs.\ $-1.33$): real humans are more price-driven than the AI simulates. The one sign flip is log historical price (human: $-0.067$, agent: $\approx 0.0$): humans slightly penalize expensive hotels while the agent is indifferent to price history.

\subsection*{Model and prompts}

\textbf{Model:} Claude Opus 4.6 Extended (Anthropic).

\textbf{Part 1: Zero-shot simulation.} The full dataset (153,009 rows, 8,354 queries) was stripped of the \texttt{booking\_bool} column. The model was prompted:

\begin{quote}
``You are simulating a customer searching for a hotel on Expedia. You will be given the customer's travel context and a list of available hotels with their attributes. Choose either one hotel to book, or choose `no purchase'. Respond in column O (\texttt{booking\_bool}) with only a single integer: 1 if you choose a hotel for a customer, or 0 for no purchase. Make sure you assign a maximum of one hotel to each customer, but it can also be none if the option is not appealing.''
\end{quote}

No real booking decisions were shown as examples. This is zero-shot simulation of customer behavior.

\textbf{Part 2: Large-context few-shot prediction.} The data was split 80/20. Both \texttt{train\_80.csv} and \texttt{test\_20.csv} were uploaded simultaneously. The prompt was extended with: \textit{``Use \texttt{train\_80.csv} as context to predict the empty \texttt{booking\_bool} column in \texttt{test\_20.csv}.''}

\textbf{Methodological limitation.} The training file contains approximately 122,400 rows. No current LLM attends uniformly to 122k CSV rows in a single context. The model applies its general prior modulated by whatever patterns it can extract from the portion of the training file within its attention window. This is a structural disadvantage relative to MLE, which integrates all 153k observations with equal weight.

\subsection*{Model's stated reasoning (zero-shot)}

The simulated customer evaluates each hotel by weighing quality against price. Quality is a composite of star rating (weighted most heavily), review score, location score, brand recognition, and accessibility, balanced against how the hotel's price compares to others in the same search and to its own historical pricing. Promotions provide an additional nudge. The customer identifies the best-value option but filters out any hotel below a baseline threshold (at least 2 stars, review score $\geq 2$).

Even when a strong option exists, the customer does not always book. The model applies a probabilistic conversion rate of roughly 4--8\%, nudged upward when the best option is especially compelling, yielding an overall 8.4\% booking rate. \textbf{Important caveat:} the model explicitly \textit{designed} its conversion rate to match platform averages. This is a deliberate engineering choice, not behavior emergent from the data, which partly explains the intercept divergence ($-5.78$ agent vs.\ $-1.75$ human).

\subsection*{Predictive evaluation}

The data was split 80/20. The AI agent received the training set as context and predicted booking decisions on the held-out test set. The MNL from Problem 1 produced choice probabilities $\mathbb{P}(j|S)$ on the same queries.

\textbf{Metrics.} Two complementary metrics are used. \textit{Session top-1 accuracy} asks whether a model selects the same hotel the real customer booked (argmax of predicted probabilities vs.\ realized choice, per session). It is a strict exact-match bar: individual choices are noisy realizations from a distribution, so even the true model will miss most. This makes top-1 useful for ranking models but a poor absolute measure of fit. \textit{Distributional behavior comparison} asks whether models reproduce aggregate booking rates across feature bins (star rating, promotion, price quartile) and the average attribute profile of booked hotels. This is the primary metric because it tests whether a model captures the statistical regularities in choice behavior -- the thing that matters for pricing and assortment decisions -- rather than predicting individual noise.

\textbf{Results.} Session top-1: MNL 29.1\% vs.\ agent 19.6\% (487 vs.\ 327 correct out of 1,671 sessions). The MNL outperforms the agent by a substantial margin on exact-match prediction.

\subsection*{Behavior comparison}

We check whether each model reproduces the \textit{aggregate behavior} in the data: booking rates across attribute bins and the average attribute profile of booked hotels.

\begin{table}[h]
\centering
\begin{tabular}{lrrrr}
\toprule
Star rating & $n$ & Real & Agent & MNL \\
\midrule
1 & 93     & 0.032 & 0.022 & 0.020 \\
2 & 6,674  & 0.029 & 0.029 & 0.031 \\
3 & 12,602 & 0.040 & 0.035 & 0.039 \\
4 & 8,545  & 0.049 & 0.055 & 0.045 \\
5 & 1,809  & 0.038 & 0.032 & 0.042 \\
\bottomrule
\end{tabular}
\caption{Booking rate by star rating.}
\end{table}

\begin{table}[h]
\centering
\begin{tabular}{lrrrr}
\toprule
Promotion flag & $n$ & Real & Agent & MNL \\
\midrule
0 (no promotion) & 25,243 & 0.036 & 0.031 & 0.036 \\
1 (promotion)    &  4,480 & 0.061 & 0.088 & 0.055 \\
\bottomrule
\end{tabular}
\caption{Booking rate by promotion flag.}
\end{table}

\begin{table}[h]
\centering
\begin{tabular}{lrrrr}
\toprule
Price quartile & $n$ & Real & Agent & MNL \\
\midrule
Q1 (low)  & 7,477 & 0.042 & 0.052 & 0.040 \\
Q2        & 7,405 & 0.049 & 0.049 & 0.044 \\
Q3        & 7,437 & 0.040 & 0.038 & 0.042 \\
Q4 (high) & 7,404 & 0.029 & 0.018 & 0.030 \\
\bottomrule
\end{tabular}
\caption{Booking rate by price quartile.}
\end{table}

\begin{table}[h]
\centering
\begin{tabular}{lrrr}
\toprule
Feature & Real bookings & Agent bookings & MNL weighted \\
\midrule
Star rating       & 3.300  & 3.332  & 3.283  \\
Review score      & 4.036  & 4.073  & 4.047  \\
Price (\$)        & 121.98 & 110.04 & 123.69 \\
Promotion rate    & 0.228  & 0.337  & 0.214  \\
Brand rate        & 0.747  & 0.779  & 0.746  \\
\bottomrule
\end{tabular}
\caption{Attribute profile of booked hotels. MNL column shows preference-weighted averages.}
\end{table}

\textbf{Results.} The MNL matches the real booking rate almost exactly in every bin, expected since it was fit to maximize likelihood on data with the same structure. The agent tracks the real pattern directionally but overreacts to promotions (8.8\% booking rate on promoted hotels vs.\ real 6.1\%) and underreacts to high prices (1.8\% of Q4-price hotels booked vs.\ real 2.9\%). The attribute profile confirms this: the agent books hotels that are approximately \$12 cheaper on average and 10 percentage points more likely to be promoted than real customers do. Despite matching the overall booking rate (3.9\% vs.\ 4.0\%), the agent systematically reallocates which hotels get booked toward the cheaper-and-promoted end of the spectrum.

\textbf{Discussion.} The agent captures the qualitative structure of human hotel choice but cannot calibrate the relative magnitudes. The MNL, fit by maximum likelihood on 153k observations, learns the exact statistical tradeoffs from revealed preferences. The agent relies on general reasoning with no access to the aggregate patterns in the data during inference. This shows up in the behavior comparison: the agent overweights promotional discounts and underweights price levels relative to real customers, even when those customers face the same choice sets. For structured choice settings like hotel search, statistical models trained on observed behavior reproduce aggregate choice patterns far more accurately than general-purpose reasoning.

%%===================================================================
\appendix
\section{Saturday vs.\ Non-Saturday Segmentation}
\label{sec:sat}
%%===================================================================

This appendix reports the Saturday/non-Saturday segmentation (the example split named in the problem statement) as a baseline against which to evaluate the solo/group choice in Problem 6.

\textbf{Mix:} 4,817 Saturday queries (57.7\%), 3,537 non-Saturday queries (42.3\%).

\begin{table}[h]
\centering
\begin{tabular}{lrrr}
\toprule
Feature & Saturday & Non-Saturday & Diff \\
\midrule
Intercept            & $-1.755$ & $-1.739$ & $+0.016$ \\
Star rating          & $+0.358$ & $+0.474$ & $+0.116$ \\
Review score         & $+0.111$ & $+0.110$ & $-0.001$ \\
Brand                & $+0.098$ & $+0.109$ & $+0.010$ \\
Location score       & $+0.026$ & $+0.017$ & $-0.009$ \\
Accessibility        & $+0.036$ & $+0.052$ & $+0.016$ \\
Log hist.\ price     & $-0.062$ & $-0.070$ & $-0.009$ \\
Price (USD)          & $-1.155$ & $-1.569$ & $-0.414$ \\
Promotion flag       & $+0.168$ & $+0.148$ & $-0.020$ \\
\bottomrule
\end{tabular}
\caption{MMNL coefficients for Saturday vs.\ non-Saturday split. Compare to Problem 6 (Table~5): differences here are far smaller.}
\end{table}

Most coefficients differ by less than 0.02. The only meaningful gap is price sensitivity ($-1.16$ vs.\ $-1.57$, difference 0.41), which is smaller than the early/late price gap (0.63, Problem 4) and less than one-third the solo/group intercept gap (0.50, Problem 6). All other features are essentially identical across the two types.

\begin{table}[h]
\centering
\begin{tabular}{lrrrrrr}
\toprule
Dataset & $|S|$ & $R_\text{mix}$ & gap$_\text{Sat}$ & gap$_\text{Non-Sat}$ \\
\midrule
data1 & 18 & 107.35 & 0.010 & 0.002 \\
data2 & 10 & 130.95 & 0.000 & 0.000 \\
data3 & 18 & 121.17 & 0.000 & 0.000 \\
data4 & 11 &  97.26 & 0.000 & 0.000 \\
\bottomrule
\end{tabular}
\caption{Assortment gaps under Saturday/non-Saturday MMNL. Compare to Problem 5 (max gap \$0.92) and Problem 6 (max gap \$3.36).}
\end{table}

All assortment gaps are at or near zero. In three of four datasets the compromise assortment is simultaneously optimal for both types. The maximum gap (\$0.01 on data1) is two orders of magnitude smaller than the solo/group gap (\$3.36 on data2). This confirms that Saturday/non-Saturday encodes no decision-level disagreement: both types want the same assortment, so there is no value in personalization along this dimension.

\end{document}
